\documentclass[10pt,a4paper]{article}
\newcommand{\ReportShortTitle}{Encoder geometry and multiuser sum design}
\usepackage[T1]{fontenc}
\usepackage[utf8]{inputenc}
\usepackage{lmodern}
\usepackage{microtype}
\usepackage[a4paper,margin=24mm,headheight=14pt]{geometry}
\usepackage{amsmath,amssymb,mathtools}
\usepackage{booktabs,tabularx,array}
\usepackage{enumitem}
\usepackage{xcolor}
\usepackage{fancyhdr}
\usepackage{titlesec}
\usepackage[hidelinks]{hyperref}

\definecolor{ink}{HTML}{18212B}
\definecolor{accent}{HTML}{135F73}
\definecolor{muted}{HTML}{5E6A73}
\definecolor{pale}{HTML}{EEF5F7}
\color{ink}
\hypersetup{colorlinks=true,linkcolor=accent,urlcolor=accent}
\setlength{\parindent}{0pt}
\setlength{\parskip}{0.55em}
\setlist{nosep,leftmargin=1.55em}
\setcounter{secnumdepth}{2}
\titleformat{\section}{\Large\bfseries\color{accent}}{\thesection}{0.6em}{}
\titleformat{\subsection}{\large\bfseries\color{ink}}{\thesubsection}{0.55em}{}
\pagestyle{fancy}
\fancyhf{}
\fancyhead[L]{\small ODMA--URA research record}
\fancyhead[R]{\small\color{muted}\ReportShortTitle}
\fancyfoot[C]{\thepage}
\renewcommand{\headrulewidth}{0.3pt}
\newcolumntype{Y}{>{\raggedright\arraybackslash}X}
\newcommand{\R}{\mathbb{R}}
\newcommand{\C}{\mathbb{C}}
\newcommand{\Z}{\mathbb{Z}}
\newcommand{\F}{\mathbb{F}}
\newcommand{\one}{\mathbf{1}}
\newcommand{\given}{\,|\,}
\newcommand{\CN}{\mathcal{CN}}
\DeclareMathOperator{\diag}{diag}
\DeclareMathOperator{\Cov}{Cov}
\DeclareMathOperator{\Var}{Var}
\DeclareMathOperator{\supp}{supp}
\DeclareMathOperator*{\argminop}{arg\,min}
\DeclareMathOperator*{\argmaxop}{arg\,max}
\newcommand{\evidence}[1]{\textcolor{accent}{\textbf{Evidence status:}} #1}
\providecommand{\ReportUpdated}{30 August 2026}
\newcommand{\ProjectReport}[4]{%
  \pagenumbering{roman}
  \begin{titlepage}
    \vspace*{18mm}
    {\large\color{accent}\textbf{ODMA--URA research record}}\\[12mm]
    {\huge\bfseries #1}\\[5mm]
    {\Large\color{muted}#2}\\[18mm]
    \begin{tabularx}{\textwidth}{@{}p{0.16\textwidth}Y@{}}
      Stage & #3\\
      Updated & \ReportUpdated\\
      Status & #4
    \end{tabularx}
    \vfill
    \textcolor{muted}{This report is a chronological research record. Measured results, interpretations, and open
    hypotheses are labelled separately. Detailed numerical ledgers remain in \texttt{results/}.}
  \end{titlepage}
  \pagenumbering{arabic}
}

\newcommand{\E}{\mathbb E}
\newcommand{\tr}{\operatorname{tr}}
\newcommand{\GF}{\operatorname{GF}}
\begin{document}
{\LARGE\bfseries Encoder geometry and multiuser sum design}\par
{\large Polar/CRC signatures, placement, energy and matched comparisons}\par
{\color{muted}30 October 2026 --- Encoder algebra; companion to the exact inference in Report 7}

\textbf{Result to explain.} With unchanged full-alphabet D2, learned symbol energies bring ODMA--polar close to
fixed dense codebooks. Full-grid mean PUPE at $B=12/14$ is $0.17116/0.28776$, versus fixed polar
$0.18346/0.29620$ and dense $0.17290/0.28526$. These are three-seed, distinct-message means, not an equivalence proof.
We derive the balancing correction, identify the structure it cannot remove, and distinguish a geometric certificate
from a decoding claim. No new receiver or production training is introduced here.

\section{The encoder and its exact second moment}
Split $w=(p,v)$ into $B_p$ placement bits and $B_c$ coded-payload bits. Put $P=2^{B_p}$, $V=2^{B_c}$ and $M=PV$.
The linear CRC/polar generator $L\in\GF(2)^{N_c\times B_c}$ produces the sign vector
\begin{equation}
 b_i(v)=(-1)^{\ell_i^T v},\qquad
 \phi(p,v)=S_pD_pb(v),\qquad D_p=\diag(\sqrt{e_p}),\quad \sum_i e_{p,i}=1.
 \label{eq:encoder}
\end{equation}
Here $\ell_i^T$ is row $i$ of $L$; $S_p\in\R^{n\times N_c}$ inserts symbols at distinct positions $r(p,i)$.
Thus $S_p^TS_p=I$, and every waveform has energy one. The learnable bank $e$ has shape $P\times N_c$,
not $N_c\times M$. The current implementation uses $e_p=\operatorname{softmax}(\theta_p)$, so energies are positive.
Polar bits and placements remain fixed. The explicit reference is
\begin{equation}
 \Phi=[S_1D_1B_{\rm pol}\ \cdots\ S_PD_PB_{\rm pol}]\in\R^{n\times M},\qquad
 B_{\rm pol}=[b(v)]_v\in\R^{N_c\times V},\qquad y=\Phi a^\star+z.
\end{equation}
The count vector $a^\star$ sums to the known activity $K$; the real noise is $z\sim\mathcal N(0,\sigma^2I)$.

For a uniform payload, a nonzero binary linear function is zero or one equally often. Consequently,
\begin{equation}
 Q_{ij}:=\E_v[b_i(v)b_j(v)]
 =\E_v[(-1)^{(\ell_i\oplus\ell_j)^Tv}]
 =\begin{cases}1,&\ell_i=\ell_j,\\0,&\ell_i\ne\ell_j.\end{cases}
 \label{eq:Q}
\end{equation}
So repeated generator rows mean \emph{identical symbol functions}, not merely correlated samples. We can compute
the whole codebook's normalized frame operator without enumerating payloads:
\begin{equation}
 \boxed{C:=\frac1M\Phi\Phi^T=\frac1P\sum_p S_pD_pQD_pS_p^T},\qquad \tr C=1.
 \label{eq:C}
\end{equation}
$C$ is an uncentered second moment, not generally a covariance. Its diagonal is particularly simple:
\begin{equation}
 C_{rr}=\frac{q_r}{P},\qquad q_r=\sum_{p,i:\,r(p,i)=r}e_{p,i}.
 \label{eq:rows}
\end{equation}
Unit energy fixes each \emph{column} of $\Phi$; balanced $q_r$ fixes total energy across its \emph{rows}.
They are different conditions. Current sizes are $n=256,N_c=32,B_c=6$, with $P=64/256$ at $B=12/14$.

\textbf{Computation and checks.} After forming $L,Q$, moments cost $O(PN_c^2+n^2)$ with no $2^{B_c}$ table;
balancing costs $O(PN_c+n)$ per iteration. The pattern bank still scales as $2^{B_p}$, so keep $B_p$ bounded when
discussing large $B$. Encoding costs $O(N_c\log N_c+n)$ with dense waveform output; D2 still keeps $2^B$ states.
\texttt{benchmarks/polar\_energy\_geometry.py} computes the diagnostics. Regression tests verify explicit-matrix
agreement, the analytic gradient, balancing constraints and the known-$K$ covariance identity below.

\clearpage
\section{Why balancing is relevant to D2}
Let $G=\Phi^T\Phi$. For unit columns, D2's scalar statistic has the exact decomposition
\begin{equation}
 t_j=\phi_j^T(y-\Phi\widehat a)+\widehat a_j
 =a_j^\star+\sum_{k\ne j}G_{jk}(a_k^\star-\widehat a_k)+\phi_j^Tz.
\end{equation}
If the other errors were zero-mean, uncorrelated and of common variance $\nu$, the interference variance would be
$\nu\sum_{k\ne j}G_{jk}^2$. D2 approximates this with the mean posterior variance; learned scaling and damping then
adjust it. Exact-$K$ dependence, unequal uncertainties and observation reuse prevent this from being exact inference.

Nevertheless, its average geometric interference factor has an exact identity:
\begin{align}
 \frac1M\sum_j\sum_{k\ne j}G_{jk}^2
 &=\frac{\|\Phi^T\Phi\|_F^2-M}{M}=M\tr(C^2)-1\\
 &=\frac Mn-1+M\left\|C-\frac{I}{n}\right\|_F^2.
 \label{eq:risk}
\end{align}
Since $\tr C=1$, the lower bound is $M/n-1$, attained exactly when $C=I/n$. This is the unit-norm tight-frame
criterion underlying the classical frame potential~\cite{frame}. It is an optimum of \emph{average squared
correlation}, not of PUPE. All remaining questions depend on that distinction.

For interpretation, use the dimensionless defect
\begin{equation}
 \boxed{\Delta:=n\left\|C-\frac{I}{n}\right\|_F^2
 =\underbrace{n\sum_r(C_{rr}-1/n)^2}_{\Delta_{\rm row}}
 +\underbrace{n\sum_{r\ne s}C_{rs}^2}_{\Delta_{\rm off}}.}
 \label{eq:split}
\end{equation}
The first term is exactly the squared coefficient of variation of row energies. The second records correlations
between physical coordinates. Balancing removes the first term; it does not automatically remove the second.
The participation measure used in the result review is $1/\tr(C^2)=n/(1+\Delta)$, not the literal matrix rank.

\textbf{A precise conditional equivalence.} If every row of $L$ is distinct, then $Q=I$ and $C$ is diagonal.
If the support admits balanced positive energies, setting $q_r=P/n$ therefore gives $C=I/n$ exactly.
The structured encoder then attains the unrestricted tight-frame minimum in \eqref{eq:risk}. This is a useful
formal statement, but it neither makes its Gram matrix equal to a dense codebook's nor makes decoding equivalent.

\textbf{Our current code does not satisfy that premise.} The effective $32\times6$ generator has 29 distinct
rows, including one zero row. For $r\ne s$, \eqref{eq:C} gives
\begin{equation}
 C_{rs}=\frac1P\sum_{p,i\ne j:\,r(p,i)=r,\,r(p,j)=s}
     \sqrt{e_{p,i}e_{p,j}}\,\mathbf1\{\ell_i=\ell_j\}\ \geq 0.
 \label{eq:duplicates}
\end{equation}
With strictly positive energies, a repeated row creates a positive off-diagonal contribution which other patterns
cannot cancel in this family. Thus exact $C=I/n$ is impossible here at finite logits. Allowing some energies to vanish
can remove individual contributions but changes the effective code and may damage distances; it is not a free fix.
This is a statement about the selected small-payload construction, not about all polar codes or large $B$.

\clearpage
\section{Deriving the deterministic correction}
Treat each allowed placement $(p,i)$ as an edge joining pattern $p$ to resource $r(p,i)$. We want
\begin{equation}
 e_{p,i}\geq0,\qquad \sum_i e_{p,i}=1,\qquad
 \sum_{p,i:\,r(p,i)=r}e_{p,i}=\frac Pn.
 \label{eq:marginals}
\end{equation}
These are linear constraints: unit mass per pattern and equal total mass per channel coordinate. Among all feasible
allocations, choose the one nearest the original uniform profile $e^0_{p,i}=1/N_c$ in generalized KL divergence:
\begin{equation}
 \min_{e\text{ satisfying }\eqref{eq:marginals}}
 \sum_{p,i}\left[e_{p,i}\log\frac{e_{p,i}}{e^0_{p,i}}-e_{p,i}+e^0_{p,i}\right].
 \label{eq:KL}
\end{equation}
This is a strictly convex problem on a convex feasible set, so its energy solution is unique when feasible.
For an interior solution the Lagrangian derivative is
\begin{equation}
 \log(e_{p,i}/e^0_{p,i})+\alpha_p+\beta_{r(p,i)}=0,
 \qquad e_{p,i}=e^0_{p,i}\,u_pv_{r(p,i)},
 \label{eq:scale}
\end{equation}
with positive scaling factors $u,v$. This motivates alternating two exact marginal corrections:
\begin{equation}
 e_{p,i}\ \leftarrow\ e_{p,i}\frac{P/n}{q_{r(p,i)}},
 \qquad
 e_{p,i}\ \leftarrow\ \frac{e_{p,i}}{\sum_j e_{p,j}}.
 \label{eq:ipf}
\end{equation}
This is iterative proportional fitting, or matrix scaling on the allowed edges~\cite{csiszar}.
Starting at $e^0$, it preserves the multiplicative form \eqref{eq:scale}. A positive converged solution satisfies
both the constraints and the stationarity equations, hence solves \eqref{eq:KL}. With a strictly positive feasible
allocation, the standard alternating-projection result applies. Boundary solutions need separate care.

\textbf{Feasibility is not automatic.} Every resource must be covered, but coverage is insufficient.
For nonnegative allocations, a necessary and sufficient transport condition is
\begin{equation}
 |\mathcal P|\leq\frac Pn\,|N(\mathcal P)|\quad\text{for every subset of patterns }\mathcal P,
\end{equation}
where $N(\mathcal P)$ is their set of reachable resources. It follows from the source-pattern-resource-sink
max-flow cut condition. Strict positivity on every allowed edge is an additional requirement.
Our implementation stops only after checking both marginals and reports nonconvergence rather than silently
labelling an allocation balanced. All six current banks converge to positive solutions.

\subsection{Going beyond balancing without guessing another encoder family}
The natural next analytic objective is the full defect $\mathcal J=\|C-I/n\|_F^2$, not only its diagonal part.
For $a_p=\sqrt{e_p}$ and $H_p=S_p^T(C-I/n)S_p$, direct differentiation of \eqref{eq:C} gives
\begin{equation}
 \frac{\partial\mathcal J}{\partial e_{p,i}}
 =\frac{2}{P\sqrt{e_{p,i}}}\sum_j(H_p)_{ij}Q_{ij}\sqrt{e_{p,j}},\qquad
 \frac{\partial\mathcal J}{\partial\theta_{p,i}}
 =e_{p,i}\left(g_{p,i}-\sum_j e_{p,j}g_{p,j}\right),
 \label{eq:gradient}
\end{equation}
where $g=\partial\mathcal J/\partial e$. The gradient has been checked against automatic differentiation.
It gives a decoder-independent diagnostic or candidate initialization/loss. Unlike \eqref{eq:KL}, this general
objective is not claimed convex, nor is minimizing it claimed to minimize PUPE. No new training arm is implemented.

\clearpage
\section{What remains after balancing: measured and unmeasured}
The following moments are evaluated exactly from the six returned energy checkpoints and original codebooks,
using \eqref{eq:C}; each entry averages three seeds. The balanced construction uses no received data or trained energies.
These are new geometry diagnostics, not new decoding experiments.

\begin{center}
\begin{tabular}{llrrr}
\toprule
$B$ & Construction & $\Delta_{\rm row}$ & $\Delta_{\rm off}$ & $\Delta$\\
\midrule
12 & Original uniform & 0.102865 & 0.024089 & 0.126953\\
   & Deterministic balance & 0 & 0.023249 & 0.023249\\
   & Learned energies & 0.003583 & 0.014004 & 0.017586\\
   & Fixed dense & 0.000526 & 0.061400 & 0.061926\\
\midrule
14 & Original uniform & 0.026324 & 0.006607 & 0.032931\\
   & Deterministic balance & 0 & 0.006567 & 0.006567\\
   & Learned energies & 0.002477 & 0.005926 & 0.008402\\
   & Fixed dense & 0.000134 & 0.015481 & 0.015615\\
\bottomrule
\end{tabular}
\end{center}

At $B=12$, learning both nearly balances rows and reduces the off-diagonal term by about 40\% relative to balancing
alone. This identifies a correction beyond inverse resource degree: energy is moved away from harmful repeated-symbol
interactions. It establishes what changed, not how much of the PUPE gain each term caused.
At $B=14$, learning reduces the off-diagonal term by about 10\%, but its residual row imbalance makes total
$\Delta$ larger than deterministic balancing's. Thus the learned objective is not simply minimizing frame potential.
More decisively, at $B=14$ both balanced and learned polar have smaller defect than the \emph{actual fixed dense}
reference, despite dense's lower measured PUPE. Reducing average interference cannot by itself explain or guarantee
closing that remaining gap. The earlier small-grid receiver swaps support an encoder benefit, but are not a
full-grid causal decomposition.

\subsection{Three reasons that a residual performance gap can survive}
\textbf{1. The codeword mean and activity law.}
Let $m_i=\mathbf1\{\ell_i=0\}$. Then
\begin{equation}
 \mu=\E_w\phi(w)=\frac1P\sum_p S_pD_pm,\qquad
 \Sigma=C-\mu\mu^T.
\end{equation}
For IID uniform messages the noiseless sum has covariance $K\Sigma$; for $K$ uniformly selected distinct messages it is
$K(M-K)/(M-1)\,\Sigma$. These identities are exact. A nonzero $\mu$ is not random information; its common contribution
$K\mu$ can be subtracted when $K$ is known. That does not recover its spent transmission energy. Iterative D2 error
covariance is different again, so neither $C$ nor $\Sigma$ alone gives its true conditional uncertainty.

\textbf{2. Rare correlations and competing message sets.}
Frame potential averages squared correlations; it can hide a damaging tail. For two allowed count vectors $a,a'$, the
actual separation is
\begin{equation}
 d^2(a,a')=\|\Phi(a-a')\|^2=(a-a')^TG(a-a').
\end{equation}
Under AWGN, the binary equal-prior comparison error is $\mathcal Q(d/(2\sigma))$. Here $\mathcal Q$ is the Gaussian
tail, unrelated to the symbol matrix $Q$. The multiplicity of close competing sets matters too. Multiuser aliases,
within-pattern weighted distances and rare cross-pattern correlations are therefore separate diagnostics.

\textbf{3. The receiver and its optimization.}
D2 replaces correlated, non-Gaussian uncertainty by calibrated scalar channels. Its loss and early stopping do not
optimize PUPE directly. A remaining gap can be waveform geometry, decoder approximation or training; it cannot be
assigned to the encoder solely from the displayed defect. Frozen-calibration swaps are informative but not optimal decoding.

\clearpage
\section{Design for competing sums, with balancing as one constraint}
\label{sec:sum-design}
The reference encoder remains fixed. The following identities concern one known pattern and a supplied
occupancy $k$, or a conditional residual after fixing the outside messages. They are not an oracle assumption
for a complete receiver. The primary source is IID uniform; the distinct-message control changes legal counts
and prior weights. The loss remains multiset-PUPE.

\subsection{The noisy statistic and the distance that energy actually changes}
For each observed character $g$, put $E_{p,g}=\sum_{i:\ell_i=g}e_{p,i}$ and
$x_g(a)=\sum_v a_v\chi_g(v)$. Independent real AWGN gives the sufficient statistics
\begin{equation}
 t_{p,g}=\frac{\sum_{i:\ell_i=g}\sqrt{e_{p,i}}\,y_{r(p,i)}}{E_{p,g}},\qquad
 t_{p,g}=x_g(a)+\nu_g,\quad \nu_g\sim\mathcal N(0,\sigma^2/E_{p,g}),
 \label{eq:character-statistic}
\end{equation}
independently across $g$. Project each group of repeated-character coordinates onto its amplitude vector
to prove sufficiency; the perpendicular observation contributes only a common likelihood factor.
Consequently the exact squared separation of two same-pattern count vectors is
\begin{equation}
 d_p^2(a,a')=\sum_g E_{p,g}\big[x_g(a)-x_g(a')\big]^2.
 \label{eq:character-distance}
\end{equation}
Positive energies preserve the integer kernel. Repeated rows add precision to existing characters;
they provide no new character equation. Splitting energy differently among identical rows cannot change
this conditional channel if their total $E_{p,g}$ is preserved. It can change mixed-pattern interactions.
The zero row measures occupancy, and carries no payload information when that occupancy is known.
Unknown pattern occupancies prevent deleting it from the full system without further analysis.

Report 7 derives $F(U,d)=\{g\in U^\perp:g^Td=1\}$ for two distinct cosets of an
$r$-dimensional payload subspace. Substituting its Fourier identity gives
\begin{equation}
 \boxed{d_p^2(a+U,a+d+U)=4^{r+1}\sum_{g\in F(U,d)} E_{p,g}.}
 \label{eq:flat-distance}
\end{equation}
Missing flats give exact aliases; low-energy flats give near competitors. This criterion is complete
for the coset family only. General integer count differences still use \eqref{eq:character-distance}.

For any two full-frame hypotheses with sums $s,s'$ and source masses $\pi,\pi'$, the exact log posterior
odds are
\begin{equation}
 \log\frac{\pi}{\pi'}+
 \frac{[y-(s+s')/2]^T(s-s')}{\sigma^2}.
 \label{eq:binary-odds}
\end{equation}
Under the first hypothesis the likelihood term has mean $d^2/(2\sigma^2)$ and variance $d^2/\sigma^2$.
Its binary MAP error is $\mathcal Q(d/(2\sigma)+(\sigma/d)\log(\pi/\pi'))$ for $d>0$.
IID count masses contain factorials; distinct legal counts have equal prior mass. Pairwise odds describe
one competition, not a full posterior or PUPE. Competing-set multiplicity and the count loss still matter.

\subsection{A finite repair obstruction and two explicitly changed signatures}
\evidence{Exhaustive certificates for the six-bit generator, not receiver performance.}
Report 7 classifies all its first four-user aliases by 58 missing three-flats. Adding character set $D$
separates every old alias exactly when $D\cap F\ne\varnothing$ for every such flat.
All $\binom{35}{6}=1,623,160$ six-element subsets of the missing characters were checked with integer
bit masks: the maximum coverage is 55 of 58. A seven-character cover is
\[
 D=\{8,13,19,52,60,61,63\}.
\]
Thus seven new characters are necessary and sufficient to separate the old aliases while retaining
the old observations. Necessity also applies to any linear-character redesign which must separate those particular
pairs: removing old rows cannot help distinguish them. This is not a lower bound of 39 on the length
of every possible six-bit $B_4$ code.

With only three additions the maximum coverage is 40, attained by $\{8,20,61\}$.
Replace the second occurrences of rows 12, 36 and 2, respectively, by these rows. All old distinct
characters remain. Splitting the old classes gives 144 ambiguous pairs, or 288 ambiguous multisets,
and minimum single-user Hamming distance stays 12. No new equality can arise because no old character
was removed. This choice optimizes the number of hit flats, not necessarily Bayes risk or noisy separation.
Its $b=6,N_c=32,n,P$, unit power and net payload rate can be kept unchanged. It changes the emitted
payload signature, not the CRC rule or placements. Ordinary CRC-SCL at the old symbol positions is
no longer a matched receiver. Unique rows yield $Q=I$, but that second-moment improvement alone is insufficient.

Alternatively, append the seven rows to the 32 reference symbols, retaining their original energy profile
scaled by $1-\eta$ and assigning energy $\eta/7$ to each new row, with $0<\eta<1$.
The resulting 39-symbol signature has no same-pattern $k\leq4$ alias: the old observations separate
all previously unequal sums, and $D$ separates every old equality. Each old alias acquires
$d^2\geq64\eta/7$; every old nonzero separation retains at least $(1-\eta)d_{\rm old}^2$.
This is a noiseless construction plus a distance bound, not a polar mother code of length 39.
It retains six payload bits and unit power but needs seven extra occupied coordinates per pattern;
at fixed $n$ those placements and their increased overlaps must be specified. Payload bits per occupied
symbol fall from $6/32$ to $6/39$. Global balancing is not automatic. The original polar symbols can
guide proposals, but the auxiliary parities require matched scoring and possibly new conditional operations.

The discovery cost here includes full small-alphabet alias enumeration and a combinatorial cover search.
Neither operation scales automatically with payload length. A fixed-length redesign that also deletes
old distinct characters could create new aliases; the seven-row cover alone would not certify it.

\subsection{A discriminating balancing check and a mixed-pattern counterexample}
On the matched B12/$n=256$/seed3201 placement bank, consider pattern 0 alone and all 2,080 two-payload
multisets. Comparing its 2,162,160 unordered pairs of sums gives
\begin{center}
\begin{tabular}{@{}lrr@{}}\toprule
Energy & Minimum squared sum distance & Comparisons shrunk from uniform\\\midrule
Uniform & 1.000000 & ---\\
Deterministic balanced & 0.801428 & 1,179,456\\\bottomrule
\end{tabular}
\end{center}
The closest balanced competitors are $\{0,43\}$ and $\{4,47\}$; 982,704 distances increase.
Under the three-row replacement above, with the same coordinates and energies, the corresponding
minimum distances are 1.500000 (uniform) and 1.271657 (balanced). This supplies a noisy-geometry
reason to retain that candidate beyond its exact-alias count. It does not compare full posterior risk,
the multiplicity of nearby competitors, or mixed-pattern minima.
Balancing converges to relative row error below $8.5\,10^{-11}$. This direct counterexample rules out
monotonic improvement of same-pattern sum distances from row balancing. It establishes neither worse
average PUPE nor a bank-wide minimum distance. The computation stores $O(V^2)$ pair states and an
$O(V^4)$ distance array; it is a finite diagnostic, not a receiver.

Same-pattern uniqueness does not extend automatically to mixed patterns. Take four length-two supports
$(0,1),(2,3),(0,2),(1,3)$, with one-bit signatures $b(v)=((-1)^v,(-1)^v)$ and uniform energy.
Each known pattern identifies every fixed-$k$ count vector. Yet the positive words on the first two
patterns sum to the same waveform as the positive words on the last two. These four distinct supports
are already globally balanced. Assign energies
\[
 (x,1-x),\ (1-x,x),\ (1-x,x),\ (x,1-x)
\]
on the four supports: global balance remains exact, but the competing sums have squared distance
$4(\sqrt x-\sqrt{1-x})^2$, zero only at $x=1/2$. Thus energy can change a mixed-pattern alias,
even though positive energy cannot change a fixed-pattern alias. In general, cross terms
$2\langle S_pD_pH z_p,S_qD_qH z_q\rangle$ can reinforce or cancel the within-pattern distances.

\subsection{Design criterion and next decision}
For fixed rows and placements, protect a declared set of nonzero same-pattern count contrasts with
constraints $\sum_i e_{p,i}(H z)_i^2\geq t$, alongside unit pattern energy, optional resource balancing
and an explicit positive energy floor. Maximizing $t$ is then a linear program. Including an exact
alias forces $t=0$, correctly exposing the need to change rows rather than energies. The contrast
discovery and separation oracle may be exponential; a finite contrast list certifies only that list.
Full mixed-pattern distances are not linear in these energies because of the amplitude cross terms.

The next discriminating comparison should retain the reference, compare its three-row replacement
under matched physical scoring, and measure source-weighted near competitors and proposal coverage.
It must include mixed patterns, fixed $b,n,K$, equal total power and charged receiver work. Report 7's
source-frequency bounds make the recorded exact four-user witnesses too rare to justify a broad
PUPE campaign on their removal alone. CRC changes, placement changes and auxiliary energy must be
separate controls. Report 11 gives a count-preserving affine operation for the reference and explains
why arbitrary sparse-Fourier measurements are not available for free.

Job038 retains the four fixed-resource arms: original/repaired signatures crossed with
uniform/balanced energy. Its primary comparison now uses the same full-alphabet D2 architecture for
all arms, adding frozen and learned dense/sparse controls from Job032 and learned symbol energies
from Job034. Each uses its matched calibration; repaired arms receive new calibration under the
same training protocol. Dense occupies all 256 coordinates, whereas sparse/polar occupy 32: channel
length and power match, occupied-resource count does not. The stable labels and complete experiment
contract are in the experiment index. No new decoder is introduced.

Historical encoder/decoder checkpoints are reevaluated under the current IID count loss. Their
float32 columns are explicitly unit-normalized in float64; the small rescaling and original training
cost are recorded. The comparison therefore does not reuse old IID membership scores or claim exact
floating-point replay. Joint historical training also optimized the encoder, whereas the repair is
fixed before D2 calibration. Common receiver architecture is a useful control, not a proof that the
Bayes-optimal encoder ordering is the same. Distinct-message training versus IID evaluation remains
common to all arms. There is no measured repair PUPE result yet.

A separate two-slot diagnostic reveals outside users and both patterns, then evaluates the exact IID
payload-pair posterior and count decision over $64^2$ ordered pairs. It includes mixed-pattern overlap
but excludes alternative pattern allocations; it is neither a complete receiver nor a full-system
Bayes risk. Native polar CRC/SCL is incompatible with the modified generator; this compatibility cost
is recorded without adding a custom receiver to the comparison.

\textbf{Prior-work boundary.} The $B_h$-code/adder-channel problem is classical
(\href{https://arxiv.org/abs/1809.07364}{Gu, 2018}). Noisy signature codes also study sum separation;
for example, \href{https://arxiv.org/abs/2206.10735}{Erdo\u{g}an, Maringer and Polyanskii, 2022}
use subset recovery and adversarial coordinate errors. Here the source admits counts, the noise is
Gaussian, and the encoder has fixed per-message energy and overlapping placements. Those guarantees
do not imply the present PUPE claims. The finite repair and balancing counterexample are the results
established here; their broader novelty remains open.

\section{Equivalence and next decisions}
\textbf{Geometric optimality is achievable as a theorem.} Equation \eqref{eq:risk} proves an average-interference
minimum; distinct symbol functions plus feasible balancing attain it. For the current repeated-symbol code, quantify
the constrained residual instead. These are precise statements worth pursuing, even if no PUPE theorem follows.

\textbf{Tightness alone cannot prove equal recovery.} Consider two unit-column $2\times4$ matrices:
\begin{equation}
 \Phi_1=[e_1,e_1,e_2,e_2],\qquad
 \Phi_2=[e_1,e_2,(e_1+e_2)/\sqrt2,(e_1-e_2)/\sqrt2].
\end{equation}
Both satisfy $\Phi\Phi^T=2I$. Yet $\Phi_1$ aliases messages even with one noiseless user, while $\Phi_2$ has four
distinct columns. Matching frame operators is much weaker than matching communication systems.

\textbf{A sufficient exact system equivalence.} If $\Phi_2=O\Phi_1\Pi$ for an orthogonal channel rotation $O$ and
a message permutation $\Pi$, white noise and exchangeable message priors give the same statistical problem after
rotating observations and relabelling messages. Equivalently, matching full Gram matrices up to that permutation is
sufficient for such an isometry on the spans. Bayes performance then agrees; D2 with matched parameters is also
equivariant to these transformations. No such relationship to the dense reference has been demonstrated.

\textbf{An empirical equivalence claim needs its own design.} Predefine an acceptable absolute PUPE margin and
operating region, then use independent seeds and paired trials to bound the performance difference. A mean close to
dense, or an interval containing zero, is not an equivalence test. Receiver inference and approximation remain fixed
in that claim; it is not global encoder optimality.

\subsection{What is settled, and a bounded encoder programme}
The selected row replacement is deterministic: use the stated binary rows and compute $Lv$ modulo 2.
For a supplied placement bank, the KL-balanced energy profile is uniquely defined when the stated
feasibility/convergence conditions hold. Seeded placement generation is reproducible, but is not a new
algebraic placement design. The finite exhaustive search that discovered the three rows must be charged;
no efficient universal row-selection rule at growing payload length has been established.

The encoder search is not exhausted. Even choosing 32 distinct characters from 64 admits
$\binom{64}{32}$ sets; arbitrary placements and larger payload dimensions add further freedom.
The practically bounded programme has three steps: test this fixed-resource repair; use failures to
identify a specific near-sum or overlap defect; then change only the responsible rows, energy constraints
or placements. The same-pattern max--min energy problem above is convex for a supplied contrast list;
finding all relevant contrasts and optimizing mixed-pattern geometry remain separate hard problems.
Second-moment balance and removal of a particular alias family do not solve these objectives.

Classical frame potential and KL projection already supply the balancing mathematics. Classical
adder-channel $B_h$ coding supplies the sum-identifiability objective. Our present contribution is the
selected-code character/alias certificate, finite repair obstruction and physical/count-aware comparison;
these are not yet a general new code construction or a performance theorem. Sparse-Walsh constructions
such as \href{https://arxiv.org/abs/1508.06336}{SPRIGHT} also motivate designed measurements, but require
measurement and receiver assumptions absent from this fixed placement/power system; see Report 11.

A useful stopping rule is empirical: if the repaired signature fails the common-D2 test, close this
finite candidate unless an independently diagnosed receiver mismatch explains the result. If it wins,
validate on new seeds and the full operating grid before broadening construction search. Matching a
learned dense comparator with a compact deterministic encoder is also valuable when its storage,
construction and occupied-resource advantages are stated. Neither outcome establishes universal optimality.

\subsection{Two active research directions}
\begin{enumerate}
\item \textbf{Encoder geometry.} Compare uniform, deterministic-balanced and learned energies with matched receivers.
      Diagnose the residual before adding a geometry loss or encoder freedom. A simple deterministic construction
      matching learned PUPE would be a positive result.
\item \textbf{Tractable polar-aware inference.} Use Report 7 to replace exponential message state, measuring loss
      against full D2 on identical waveforms. Unknown activity, pattern association, same-pattern collisions and joint
      interference remain open. Keep known-$K$ and unknown-$K$ results separate.
\end{enumerate}
The shared mixer is retired. Hash/coordinate-amplitude and sectioned-tree routes are parked as historical evidence,
not declared mathematically inferior. Polar energy has the stronger current performance/structure case, but no
completed scalable multiuser decoder yet. Broad new encoder families are unnecessary unless a diagnosed residual
points to a specific missing freedom.

\textbf{Scope of a possible conclusion.} ``Simple structured encoding plus geometric optimization matches the dense
reference under the tested receiver and conditions'' would be well supported by a successful controlled test.
``No subtle code structure can outperform well-behaved dense encoding'' is much stronger and does not follow from
second moments, D2 comparisons or these finite experiments. Dense is a useful reference, not a proven optimum.

\begin{thebibliography}{2}
\bibitem{frame} D. G. Mixon, T. Needham, C. Shonkwiler and S. Villar,
\emph{Three proofs of the Benedetto--Fickus theorem}, 2021/2023.
\href{https://arxiv.org/abs/2112.02916}{arXiv:2112.02916}. Background for the frame-potential criterion, not a URA theorem.
\bibitem{csiszar} I. Csisz\'ar, \emph{I-divergence geometry of probability distributions and minimization problems},
Annals of Probability 3(1), 146--158, 1975.
\href{https://web.mit.edu/6.291/www/handouts/csiszar_Idiv.pdf}{Author paper}. Background for KL projection and iterative fitting.
\end{thebibliography}
\end{document}
